% ---------------------------------------------------------------------------
% DRAFT SUBSECTION: What the reasoning traces say
% Source of all numbers: analysis/build_trace_features.py
%   -> results/traces/{trace_features.csv, ols_results.txt, traces_summary.md}
% Intended location: end of Section 6 (after sec:ranking-scaffolds), or as the
% bridge into the discussion. Uses only macros already defined in the master
% (\TODO). Labels are new; cross-references to existing labels are guarded
% with comments.
% ---------------------------------------------------------------------------
\subsection{What the Traces Say: Stated Reasoning and Realized Bids}
\label{sec:traces}

Every bid in our experiments is accompanied by a short free-text plan in
which the agent explains its choice (median 68 words). We analyze all
21{,}990 plans from the sealed-bid and clock cells with a transparent
dictionary of 17 keyword features---stated intent to shade below or bid
above value, rehearsal of the second-price payment rule, dominance language,
worst-case and profit-margin rhetoric, and explicit mechanism confusion
(calling the auction ``first-price''), among others (full regular
expressions and per-cell prevalence in the replication package,
\path{analysis/build_trace_features.py}).%
\footnote{The plans are stated strategies, not full chains of thought; all
claims in this subsection concern what models \emph{say}, with
Table~\ref{tab:trace-dissociation} showing that what they say tracks what
they do.}

Three regularities stand out. First, the modal plan applies
\emph{first-price} logic to the second-price auction: 53\% of sealed-bid
second-price traces state an intent to bid below value, 55\% worry about
``overpaying,'' and 31\% invoke a ``profit margin'' or ``buffer,'' while
only 2\% use dominance language and only 4\% correctly rehearse that the
winner pays the \emph{second}-highest bid. Claude-3.5-Haiku explicitly
mislabels the mechanism ``first-price'' in 10\% of its traces; no other
model does so more than 0.1\% of the time.

Second, the traces are not cheap talk: stated intent predicts the realized
bid direction almost perfectly. Among traces stating an intent to shade,
97\% of bids fall below value (mean deviation $-\$1.84$); among traces
stating an intent to bid above value, 88\% fall above (mean $+\$1.88$); the
largest errors come from the 4{,}541 traces that state no value-anchored
intent at all (mean deviation $-\$4.84$), whose plans anchor on
uninformative quantities such as ``slightly above half my value.'' In a
pooled regression of $|b-v|$ on the dictionary features with model and
lever-family fixed effects (cluster-robust by run), articulating \emph{any}
value-anchored plan---shading, overbidding, or truthful---is associated with
\$1.5--\$2.2 smaller absolute deviations than an unanchored plan; the text
features add $0.075$ of incremental $R^2$ over the fixed effects
($0.251$ vs.\ $0.176$). The traces are informative about \emph{how} an
agent errs, but stated reasoning is far from a sufficient statistic for the
error itself.
\TODO{Trace option 4 (stated-vs-revealed): parse the dollar amounts quoted
inside plans and report a calibration curve of stated vs.\ realized shade
per model; GPT-4o's overbid-intent traces execute the stated direction only
61\% of the time vs.\ Claude's 98\%.}

Third, and most important for interpreting the ranking, the levers of
Section~\ref{sec:ranking-scaffolds} produce a \emph{double dissociation}
between stated reasoning and behavior
(Table~\ref{tab:trace-dissociation}). \emph{Payoff Safety} (B2) moves bids
substantially while leaving stated reasoning untouched: the invariant it
asserts (``your bid determines \emph{if} you win, not \emph{what} you
pay'') is echoed in zero of its 600 treated traces, and payment-rule
rehearsal is flat relative to baseline ($2.8\%$ vs.\ $3.5\%$, $p=0.66$);
the plans keep the same shading rhetoric and simply shrink the shade to
pennies (``I will bid \$31.9 to avoid overbidding,'' value \$32). The
\emph{worst-case} scaffold (C1) does the opposite: it massively shifts
stated reasoning (worst-case language rises from $2\%$ to $43\%$) while
bids barely move. The \emph{Payoff Tree} moves both, and the \emph{menu}
restatement moves neither, serving as the placebo row. A classical
mediation analysis therefore fails at the first stage: whatever Payoff
Safety does, it does not work by making agents verbalize why truth-telling
is safe. Stated comprehension and behavioral improvement are separable
outcomes, and our effective descriptions improve the latter without the
former.
\TODO{Trace option 2 (formal mediation): upgrade this paragraph with the
full manipulation-check matrix over all levers and an explicit
first-stage-zero mediation bound; all inputs are in
\path{results/traces/trace_features.csv}.}
\TODO{Trace option 1 (LLM-judge taxonomy): replace/augment the dictionary
with a judged strategy taxonomy (truthful-dominance, shading-for-margin,
opponent-anchored, mechanism-confused, \ldots) validated against a
100-trace human-labeled set; needed before making the ``zero dominance
recognition'' claim load-bearing.}

Finally, the traces speak to \emph{why} comprehension fails. In the GPT-4o
first-price and third-price variants of the same grid, the model deploys a
nearly identical shading script---shading language in 70\%/79\%/72\% of
second-/first-/third-price traces and mean bids of $0.88$/$0.87$/$0.84$
of value---even though the optimum differs sharply across formats (truthful
in second-price; roughly $\tfrac{2}{3}v$ in our three-bidder first-price
cell; above value in third-price). The same words justify underbidding
where shading is optimal and where it is an error: the reasoning script is
a property of the model, not of the mechanism it faces.
\TODO{Trace option 3: the cross-mechanism script-invariance result is
currently GPT-4o--only (V12 recovered logs); run FPSB/TPSB cells for
Claude, Gemini, and Gemma with the frozen dictionary before claiming
generality.}

\begin{table}[htbp]
\centering
\small
\begin{tabular}{lcccc}
\toprule
& \multicolumn{2}{c}{Stated reasoning} & \multicolumn{2}{c}{Behavior} \\
\cmidrule(lr){2-3}\cmidrule(lr){4-5}
Lever (family) & target language & $\Delta$ vs.\ baseline &
mean $|b-v|$ & baseline \\
\midrule
Payoff Safety (B2) & payment-rule rehearsal & $3.5\%\to 2.8\%$ &
$1.95$ & $3.27$ \\
Worst-case scaffold (C1) & worst-case rhetoric & $2\%\to 43\%$ &
$3.07$ & $3.24$ \\
Payoff Tree (C1) & rule rehearsal & $7.5\%\to 26.7\%$ &
$1.29$ & $3.41$ \\
Menu restatement (B1) & --- & no shift & $3.55$ & $3.11$ \\
\bottomrule
\end{tabular}
\caption{Double dissociation between stated reasoning and bidding behavior
(pooled over four models, sealed-bid second-price cells; baselines are the
axis-specific baseline cells, cluster-robust tests in
\texttt{results/traces/ols\_results.txt}). Payoff Safety improves bids without
changing what agents say; the worst-case scaffold changes what agents say
without improving bids.
\label{tab:trace-dissociation}}
\end{table}
\TODO{Table~\ref{tab:trace-dissociation}: (i) decide whether the Payoff
Safety baseline column should quote the \texttt{spsb} cell ($3.11$/$-2.67$,
as in Section~\ref{sec:ranking-descriptions}) or the pooled axis baselines
($3.27$) --- both are in \path{results/traces/traces_summary.md};
(ii) reconcile the Claude Payoff-Safety sign ($+0.30$ in the logs vs.\
$-0.30$ in the current draft); (iii) caveat for the menu row: the menu
prompt replaces the ``second-highest'' vocabulary entirely (rule-rehearsal
prevalence is $0\%$ by construction), so the language column measures
prompt echo there, not comprehension --- consider dropping the language
entry for B1 or noting the confound in the caption.}
% ---------------------------------------------------------------------------
